Inter-annotator agreement results indicate that agreement is not uniform across annotation components. At the node-label level, unitizing Krippendorff's $\alpha_u$ is moderate for \textit{Analysis} (0.567), \textit{Procedural History} (0.454), and \textit{Rule} (0.613), and high for \textit{Background Facts} (0.879) and \textit{Conclusion} (0.724) in the span view. The sentence view produces nearly identical values (\textit{Analysis}: 0.574, \textit{Background Facts}: 0.882, \textit{Conclusion}: 0.724, \textit{Procedural History}: 0.456, \textit{Rule}: 0.610). This supports a narrower claim than the draft: mapping spans to sentences has negligible effect on $\alpha_u$, but only for the unitizing measure.

The same conclusion does not hold for overlap-based F1. Under edit-distance pairing, sentence-view F1 is substantially higher than span-view F1 for \textit{Analysis} (0.698 vs.\ 0.456), \textit{Background Facts} (0.915 vs.\ 0.697), and \textit{Rule} (0.705 vs.\ 0.505), which indicates that boundary variation is a major source of disagreement for these labels. At the span level, semantic pairing also raises F1 relative to edit-distance pairing across all labels (for example, \textit{Analysis}: 0.564 vs.\ 0.456; \textit{Rule}: 0.641 vs.\ 0.505). These higher scores should not be interpreted as evidence that semantic pairing is intrinsically the better agreement measure. Rather, semantic pairing is more permissive and absorbs some disagreement that edit-distance matching leaves exposed. For that reason, it is better presented as a sensitivity analysis or an upper-bound matching condition than as the primary agreement estimate.

Agreement on implicit intermediate-conclusion insertion is driven mainly by agreement on \emph{absence}. With edit-distance pairing, observed agreement is 0.645, $\kappa=-0.044$, positive agreement is 0.154, and negative agreement is 0.776. Semantic pairing yields nearly identical values (observed agreement 0.641, $\kappa=-0.048$, positive agreement 0.148, negative agreement 0.772). The stable pattern across both pairings is that annotators usually agree when no implicit insertion is needed, but they do not agree reliably on when an insertion should be introduced. This means the insertion decision is substantially noisier than the decision not to insert.

Direct-edge agreement is the weakest part of the annotation. As Table~\ref{tab:iaa_edge_summary} shows, observed agreement is only 0.074 under edit-distance pairing and 0.072 under semantic pairing, with strongly negative $\kappa$ values in both cases ($-0.581$ and $-0.593$, respectively). Because the aggregate result barely changes across pair-matching backends, weak direct-edge agreement does not appear to be primarily a lexical-matching artifact. The notebook outputs also show large annotator differences in edge density (edit pairing positive-rate imbalance: 0.747 vs.\ 0.326; semantic pairing: 0.742 vs.\ 0.330). Taken together, these results indicate that direct-edge decisions are unstable under both pairing strategies and should be interpreted more cautiously than node-label agreements.

Taken together, these results show that node-label agreement is substantially stronger than agreement on graph edges or implicit insertions. The latter two components remain much noisier and therefore require more cautious interpretation.
